% GENERATED FILE. Edit canonical lessons, then run npm run generate:books.
% canonical-source-hash: fnv1a64:af627af25a1d9ec5
% canonical-lessons: AR-C399-illa-idha, AR-C399-fi-hal, AR-C399-bi-sharti-an, AR-C399-talama, AR-C399-wa-illa

\chapter{If and unless}
\label{ch:ar-if-and-unless}
\hlchaptermodality{399}

\begin{chapteropening}
\textbf{By the end of this chapter:} \emph{I can say unless, in case, provided that, as long as and otherwise.}

Complete the last lesson of chapter 399: I can say unless, in case, provided that, as long as and otherwise.
\end{chapteropening}

\section[\texorpdfstring{\ar{إلا} \ar{إذا} (illā idhā)}{إلا إذا (illā idhā)}]{\ar{إلا} \ar{إذا} (illā idhā) — unless}

\label{lesson:AR-C399-illa-idha}



\begin{quote}
Before the new one: say the Arabic for literacy, then the Arabic for attendance.
\end{quote}

\subsection*{You'll want to know: \ar{إلا} \ar{إذا}}
\textbf{\ar{إلا} \ar{إذا}} — \emph{illā idhā} — \textquotedblleft{}unless\textquotedblright{}.

\begin{quote}
\ar{لن} \ar{نخرج} \ar{إلا} \ar{إذا} \ar{توقف} \ar{المطر}.
\end{quote}

\begin{quote}
\emph{lan nakhruja illā idhā tawaqqafa al-maṭar.}
\end{quote}

\begin{quote}
We will not go out unless the rain stops.
\end{quote}

After \ar{إذا} the verb is usually in the past, even for the future.

\begin{cousinweb}
\textbf{\ar{إلا} \ar{إذا}} (\emph{illā idhā}) joins the native Arabic particles \emph{illā}, \textquotedblleft{}except\textquotedblright{}, and \emph{idhā}, \textquotedblleft{}if\textquotedblright{}: except if, unless. The deeper origin of such particles is not settled.
\end{cousinweb}

\begin{grammarlens}[title={What you've built}]
One more word for saying what happens if.
\end{grammarlens}

\subsection*{Your turn}
\begin{itemize}
\raggedright
  \item \emph{Say it:} \emph{illā idhā}
  \item \emph{Say it:} \emph{illā idhā}, once more
  \item \emph{Say it:} \emph{lan nakhruja illā idhā tawaqqafa al-maṭar.} — We will not go out unless the rain stops.
  \item \emph{Recall:} say the Arabic for literacy, then the Arabic for attendance, then say \emph{illā idhā} again
\end{itemize}

\begin{tcolorbox}[breakable,colback=teal!4,colframe=teal!35!black,title={Before you move on}]
What does \ar{إلا} \ar{إذا} mean? (\textquotedblleft{}Unless\textquotedblright{}.) Say it once more.
\end{tcolorbox}
\section[\texorpdfstring{\ar{في} \ar{حال} (fī ḥāl)}{في حال (fī ḥāl)}]{\ar{في} \ar{حال} (fī ḥāl) — in case}

\label{lesson:AR-C399-fi-hal}



\begin{quote}
Before the new one: say the Arabic for attendance, then the Arabic for unless.
\end{quote}

\subsection*{You'll want to know: \ar{في} \ar{حال}}
\textbf{\ar{في} \ar{حال}} — \emph{fī ḥāl} — \textquotedblleft{}in case\textquotedblright{}.

\begin{quote}
\ar{في} \ar{حال} \ar{تأخر} \ar{القطار،} \ar{سأتصل} \ar{بك}.
\end{quote}

\begin{quote}
\emph{fī ḥāl taʾakhkhara al-qiṭār, sa-attaṣilu bik.}
\end{quote}

\begin{quote}
In case the train is late, I will call you.
\end{quote}

Also \ar{في} \ar{حالة}, \emph{fī ḥālat}: in the event of.

\begin{cousinweb}
\textbf{\ar{في} \ar{حال}} (\emph{fī ḥāl}) joins \emph{fī}, \textquotedblleft{}in\textquotedblright{}, with \emph{ḥāl}, \textquotedblleft{}a state, a condition\textquotedblright{}, built on the native Arabic root \emph{ḥ-w-l}, \textquotedblleft{}to change, to turn\textquotedblright{} (Lane). Literally, in the state that.
\end{cousinweb}

\begin{grammarlens}[title={What you've built}]
One more word for saying what happens if.
\end{grammarlens}

\subsection*{Your turn}
\begin{itemize}
\raggedright
  \item \emph{Say it:} \emph{fī ḥāl}
  \item \emph{Say it:} \emph{fī ḥāl}, once more
  \item \emph{Say it:} \emph{fī ḥāl taʾakhkhara al-qiṭār, sa-attaṣilu bik.} — In case the train is late, I will call you.
  \item \emph{Recall:} say the Arabic for attendance, then the Arabic for unless, then say \emph{fī ḥāl} again
\end{itemize}

\begin{tcolorbox}[breakable,colback=teal!4,colframe=teal!35!black,title={Before you move on}]
What does \ar{في} \ar{حال} mean? (\textquotedblleft{}In case\textquotedblright{}.) Say it once more.
\end{tcolorbox}
\section[\texorpdfstring{\ar{بشرط} \ar{أن} (bi-sharṭi an)}{بشرط أن (bi-sharṭi an)}]{\ar{بشرط} \ar{أن} (bi-sharṭi an) — provided that}

\label{lesson:AR-C399-bi-sharti-an}



\begin{quote}
Before the new one: say the Arabic for unless, then the Arabic for in case.
\end{quote}

\subsection*{You'll want to know: \ar{بشرط} \ar{أن}}
\textbf{\ar{بشرط} \ar{أن}} — \emph{bi-sharṭi an} — \textquotedblleft{}provided that\textquotedblright{}.

\begin{quote}
\ar{سأساعدك} \ar{بشرط} \ar{أن} \ar{تساعدني}.
\end{quote}

\begin{quote}
\emph{sa-usāʿiduka bi-sharṭi an tusāʿidanī.}
\end{quote}

\begin{quote}
I will help you provided that you help me.
\end{quote}

The verb after \ar{أن} takes the subjunctive: \ar{أن} \ar{تساعدني}.

\begin{cousinweb}
\textbf{\ar{بشرط} \ar{أن}} (\emph{bi-sharṭi an}) joins \emph{bi-}, \textquotedblleft{}with\textquotedblright{}, \emph{sharṭ}, \textquotedblleft{}a condition\textquotedblright{}, built on the native Arabic root \emph{sh-r-ṭ}, \textquotedblleft{}to set a mark, to stipulate\textquotedblright{} (Lane), and \emph{an}, \textquotedblleft{}that\textquotedblright{}: on condition that.
\end{cousinweb}

\begin{grammarlens}[title={What you've built}]
One more word for saying what happens if.
\end{grammarlens}

\subsection*{Your turn}
\begin{itemize}
\raggedright
  \item \emph{Say it:} \emph{bi-sharṭi an}
  \item \emph{Say it:} \emph{bi-sharṭi an}, once more
  \item \emph{Say it:} \emph{sa-usāʿiduka bi-sharṭi an tusāʿidanī.} — I will help you provided that you help me.
  \item \emph{Recall:} say the Arabic for unless, then the Arabic for in case, then say \emph{bi-sharṭi an} again
\end{itemize}

\begin{tcolorbox}[breakable,colback=teal!4,colframe=teal!35!black,title={Before you move on}]
What does \ar{بشرط} \ar{أن} mean? (\textquotedblleft{}Provided that\textquotedblright{}.) Say it once more.
\end{tcolorbox}
\section[\texorpdfstring{\ar{طالما} (ṭālamā)}{طالما (ṭālamā)}]{\ar{طالما} (ṭālamā) — as long as}

\label{lesson:AR-C399-talama}



\begin{quote}
Before the new one: say the Arabic for in case, then the Arabic for provided that.
\end{quote}

\subsection*{You'll want to know: \ar{طالما}}
\textbf{\ar{طالما}} — \emph{ṭālamā} — \textquotedblleft{}as long as\textquotedblright{}.

\begin{quote}
\ar{سأبقى} \ar{هنا} \ar{طالما} \ar{أنت} \ar{هنا}.
\end{quote}

\begin{quote}
\emph{sa-abqā hunā ṭālamā anta hunā.}
\end{quote}

\begin{quote}
I will stay here as long as you are here.
\end{quote}

With the past it also means often: \ar{طالما} \ar{حلمت} \ar{بذلك}, I often dreamed of that.

\begin{cousinweb}
\textbf{\ar{طالما}} (\emph{ṭālamā}) joins \emph{ṭāla}, \textquotedblleft{}it was long\textquotedblright{}, built on the native Arabic root \emph{ṭ-w-l}, \textquotedblleft{}to be long\textquotedblright{} (Lane), with \emph{mā}: for as long as.
\end{cousinweb}

\begin{grammarlens}[title={What you've built}]
One more word for saying what happens if.
\end{grammarlens}

\subsection*{Your turn}
\begin{itemize}
\raggedright
  \item \emph{Say it:} \emph{ṭālamā}
  \item \emph{Say it:} \emph{ṭālamā}, once more
  \item \emph{Say it:} \emph{sa-abqā hunā ṭālamā anta hunā.} — I will stay here as long as you are here.
  \item \emph{Recall:} say the Arabic for in case, then the Arabic for provided that, then say \emph{ṭālamā} again
\end{itemize}

\begin{tcolorbox}[breakable,colback=teal!4,colframe=teal!35!black,title={Before you move on}]
What does \ar{طالما} mean? (\textquotedblleft{}As long as\textquotedblright{}.) Say it once more.
\end{tcolorbox}
\section[\texorpdfstring{\ar{وإلا} (wa-illā)}{وإلا (wa-illā)}]{\ar{وإلا} (wa-illā) — otherwise}

\label{lesson:AR-C399-wa-illa}



\begin{quote}
Before the new one: say the Arabic for provided that, then the Arabic for as long as.
\end{quote}

\subsection*{You'll want to know: \ar{وإلا}}
\textbf{\ar{وإلا}} — \emph{wa-illā} — \textquotedblleft{}otherwise\textquotedblright{}.

\begin{quote}
\ar{يجب} \ar{أن} \ar{يدرس،} \ar{وإلا} \ar{سيرسب}.
\end{quote}

\begin{quote}
\emph{yajibu an yadrusa, wa-illā sa-yarsub.}
\end{quote}

\begin{quote}
He has to study, otherwise he will fail.
\end{quote}

\begin{cousinweb}
\textbf{\ar{وإلا}} (\emph{wa-illā}) joins the native Arabic \emph{wa-}, \textquotedblleft{}and\textquotedblright{}, with \emph{illā}, which here is \emph{in}, \textquotedblleft{}if\textquotedblright{}, fused with \emph{lā}, \textquotedblleft{}not\textquotedblright{}: and if not, otherwise. The deeper origin of such particles is not settled.
\end{cousinweb}

\begin{grammarlens}[title={What you've built}]
One more word for saying what happens if.
\end{grammarlens}

\subsection*{Your turn}
\begin{itemize}
\raggedright
  \item \emph{Say it:} \emph{wa-illā}
  \item \emph{Say it:} \emph{wa-illā}, once more
  \item \emph{Say it:} \emph{yajibu an yadrusa, wa-illā sa-yarsub.} — He has to study, otherwise he will fail.
  \item \emph{Recall:} say the Arabic for provided that, then the Arabic for as long as, then say \emph{wa-illā} again
\end{itemize}

\begin{tcolorbox}[breakable,colback=teal!4,colframe=teal!35!black,title={Before you move on}]
What does \ar{وإلا} mean? (\textquotedblleft{}Otherwise\textquotedblright{}.) Say it once more.
\end{tcolorbox}
